%\section*{Devoir maison}
\newpage
\begin{Exercise}[title={Pyramide}]

On veut construire une pyramide à base carrée avec des billes : une seule bille au sommet (étage 1), puis un carré de billes à chaque étage, en descendant. La pyramide visée a 7 étages.
\begin{center}
        
\end{center}
\Question Combien ai-je de billes à l'étage 2 ? 3 ? 4 ? 5 ? en partant du haut. En déduire une règle générale.\\
\emptybox{2cm}

\Question Donner l'équation de la fonction f donnant le nombre de billes aux $i^eme$ étage.
$$f(i)=..................$$

\Question On a 200 billes. Combien d'étages au maximum aura sa pyramide ? Justifiez !\\
\emptybox{2cm}



\begin{lstlisting}[style=pythonstyle, language=python]
nb_billes = 200
etage = 0
while nb_billes-(etage+1)**2 >= 0  :
    etage += 1
    nb_billes = nb_billes - etage**2
    
print(f"Pour {nb_billes} billes la pyramide fait {etage}étages.")
\end{lstlisting}
\Question Recopier et completer le tableau d'avancement des variables du programme ci-dessus.

\begin{center}
\large
\begin{tabular}{c|c|c|c|}
\textbf{Ligne} & \textbf{base} & \textbf{nb\_billes} & \textbf{etage}  \\
\hline
 &  &  & \\
 &  &  & \\
 &  &  & \\
 &  &  & \\
 &  &  & \\
  &  &  & \\
   &  &  & \\
    &  &  & \\
     &  &  & \\
      &  &  & \\
\end{tabular}
\end{center}

\Question Écrire un programme en Python affichant le nombre de billes d'une pyramide à partir de son nombre d'étage saisi par l'utilisateur. 



\end{Exercise}